\section{Experimental Design}
\label{sec:experiment}

\subsection{Research questions}

The experiment addresses three questions:

\begin{enumerate}
\item \textbf{RQ1:} To what extent do the tested DOCX-to-PDF conversion pipelines preserve accessibility information present in controlled source documents?
\item \textbf{RQ2:} Which tested accessibility properties are verified as preserved, observed as changed or lost, or unresolved under the available cross-format evidence?
\item \textbf{RQ3:} How do preservation outcomes differ between the tested LibreOffice and Google Docs conversion pipelines?
\end{enumerate}

Validator comparison is not a fourth core question. The project did not obtain a systematic PAC, Acrobat Accessibility Checker, or veraPDF dataset for all 26 outputs, so destination-only validator comparison is future work rather than a central result.

\subsection{Benchmark composition}

Table~\ref{tab:fixtures} summarizes the 13 atomic fixture contracts. Each fixture contributes one feature-level outcome per pipeline, yielding 13 $\times$ 2 = 26 cases. The fixtures are intentionally controlled rather than randomly sampled. F13 was reserved for reading order but deferred because reliable automated source-to-PDF equivalence could not be established; the frozen corpus therefore uses F01--F12 and F14 without renumbering. The table's contract descriptions state what is compared; they do not imply that the feature families have equal complexity or equal user impact.

\begin{table*}[t]
\centering\small
\caption{Frozen fixture contracts.}
\label{tab:fixtures}
\begin{tabular}{p{0.14\textwidth}p{0.17\textwidth}p{0.56\textwidth}}
\hline
Fixture & Property & Contract focus \\
\hline
F01 & Headings & Heading levels and order \\
F02 & Alternative text & Author-written image description \\
F03 & Lists & Ordered/unordered type, nesting, and list structure \\
F04 & Table & Table, row/cell, and header structure \\
F05 & Document language & Primary document language value \\
F06 & Inline language & Run-level language override within an en-US document \\
F07 & Decorative image & Meaningful versus explicitly decorative image state \\
F08 & Hyperlinks & Link existence, visible text, and URI destination \\
F09 & Document title & Visible title and document metadata identity \\
F10 & Complex table & Multi-level headers, spans, and associations \\
F11 & Footnotes & Note body and reference-to-note association \\
F12 & Equation & Native equation representation and visible expression \\
F14 & Caption & Figure/caption text and association \\
\hline
\end{tabular}
\end{table*}

\subsection{Converter conditions and provenance}

The LibreOffice workflow was deterministic and local. The Google Docs workflow was authenticated and browser-mediated, so its service backend was not independently version-pinned. The experiment date, browser version, observed PDF producer, source hashes, output hashes, file sizes, and validation records were retained in the automation log. All outputs were one-page PDFs in this corpus and passed the output-validation checks.

The experiment's primary unit is one fixture/converter pair. The 13 fixtures are controlled atomic cases, not a representative sample; features have unequal complexity and are not repeated instances of a population. Counts and feature-level matrices are therefore primary, while percentages are not used as headline results. No inferential statistics, significance tests, prevalence estimates, or weighted overall converter score are justified. A separate secondary sanity check uses three realistic multi-feature documents and the same existing contracts; its property-level observations are not merged into the primary denominator.

\subsection{Outcome accounting}

The primary evidence-aware set contains 11 \verifiedpreserved{}, 3 \observedpartial{}, 3 \altered{}, 2 \confirmedlost{}, 6 \unresolved{}, and 1 \measured{} case. Thus 26 cases were attempted and successfully converted, with one case remaining a measurement error. The historical V1 set contained 11 PRESERVED, 9 PARTIALLY\_PRESERVED, 3 ALTERED, 2 LOST, and 1 MEASUREMENT\_ERROR; the V1 record is retained in the interpretation changelog and is not silently overwritten. No percentages are used because the controlled feature units have unequal complexity and do not support prevalence interpretation.

\subsection{Integrated-document sanity check}

The secondary sanity check evaluated three realistic two-page documents: I01, a community access report; I02, an educational handout; and I03, a policy note. They used only existing contract families, and an independent raw-OOXML oracle passed 21/21 integrated source assertions. The same two pipelines produced six tagged PDFs, and the protocol generated 42 property--pipeline observations without requiring a document-level score. The secondary result distribution was 16 VERIFIED\_PRESERVED, 6 OBSERVED\_PARTIAL, 5 ALTERED, 3 CONFIRMED\_LOST, 11 UNRESOLVED\_EQUIVALENCE, and 1 MEASUREMENT\_ERROR. These observations demonstrate protocol applicability when properties coexist; they are not a prevalence estimate and are excluded from the primary 26-case denominator.


